\chapter{RL Circuits}
\label{ch:rl}
Inductors appear in RF chokes, tuned circuits, filters, relays, transformers, and
antenna loading coils. An RL circuit is the current-domain sibling of the RC
circuit: it is the same \emph{first-order system} of \cref{ch:linsys}, but now the
natural state is the inductor current \(i_L\), which holds the stored magnetic
energy.

\section{The Circuit}
A source \(u(t)\) drives a resistor \(R\) in series with an inductor \(L\)
(\cref{fig:rl-schematic}). Taking the output across the resistor gives a low-pass
response; across the inductor, a high-pass response.

\begin{figure}[htbp]
\centering
\begin{circuitikz}[scale=1]
\draw (0,0) to[sV,l=\(u(t)\)] (0,2)
      to[L,l=\(L\),i>=\(i(t)\)] (3,2)
      to[R,l=\(R\),v=\(v_R\)] (3,0) -- (0,0);
\draw (3,2) to[short,-o] (4,2) node[right]{$+$};
\draw (3,0) to[short,-o] (4,0) node[right]{$-$};
\node at (4.35,1) {\(v_o\)};
\end{circuitikz}
\caption{Series RL circuit with the output across the resistor (low-pass in the
current sense). The state variable is the inductor current \(i(t)\).}
\label{fig:rl-schematic}
\end{figure}

\section{From Schematic to Differential Equation}
KVL gives the source as the sum of the inductor and resistor voltages, and the
inductor law relates its voltage to its rate of current change:
\[
u(t)=L\frac{\dd i}{\dd t}+R\,i(t)
\qquad\Longleftrightarrow\qquad
\boxed{\,\frac{L}{R}\frac{\dd i}{\dd t}+i=\frac{1}{R}\,u(t).\,}
\]

\section{State-Space Form and the Generic Match}
With the single state \(x=i\),
\[
\dot{x}=-\frac{R}{L}\,x+\frac{1}{L}\,u .
\]
Comparing with \(\tau\dot y+y=Ku\) from \cref{ch:linsys},
\[
\tau=\frac{L}{R},
\qquad
K=\frac{1}{R},
\qquad
s_{\text{pole}}=-\frac{R}{L}=-\frac{1}{\tau}.
\]
It is the same first-order system as the RC circuit, with the time constant
\(L/R\) in place of \(RC\). Note the duality: RC stores energy in a voltage and
its time constant grows with \(R\); RL stores energy in a current and its time
constant \emph{shrinks} with \(R\).

\section{Transient Response}
For a DC step of height \(E\), and for decay from an initial current \(I_0\),
\[
i(t)=\frac{E}{R}\left(1-\ee^{-Rt/L}\right),
\qquad
i(t)=I_0\,\ee^{-Rt/L}.
\]
This is again the first-order response of \cref{ch:linsys}, read as an inductor
current (\cref{fig:rl-transient}): after one time constant \(\tau=L/R\) the
current has reached \(63.2\%\) of its final value. Inductor current cannot jump
under finite voltage---the dual of the capacitor's continuous voltage.

\begin{figure}[htbp]
\centering
\includegraphics[width=0.82\textwidth]{rl_transient}
\caption{RL inductor-current transient. Rise and decay share the time constant
\(\tau=L/R\); one \(\tau\) covers \(63.2\%\) of the change. The same shape as the
RC voltage transient---the RC/RL duality.}
\label{fig:rl-transient}
\end{figure}

\section{Frequency Response and the Root Locus}
The inductor impedance is \(Z_L=\jj\omega L\), which grows with frequency, so an
inductor opposes rapid current changes and looks like a short at DC. Output across
\(R\) is low-pass (\cref{fig:bode1}); across \(L\) it is high-pass
(\cref{fig:rc-highpass}), with corner \(\omega_c=R/L\). The single pole
\(-R/L\) rides the negative real axis (\cref{fig:rc-rootlocus}); as \(R\)
increases the pole moves \emph{left} (faster), the mirror of the RC case.

\begin{controlsbox}
\textbf{Connection.} RC and RL are the two first-order lags of circuit theory.
Same pole structure, same Bode shape, same step response---only the state
variable (voltage vs current) and the direction of the \(R\)-dependence differ.
Learn one and you have both.
\end{controlsbox}

\section{Energy Interpretation}
The inductor stores \(E_L=\tfrac12 Li^2\) in its magnetic field. When current is
interrupted, that energy must go somewhere, which is why breaking an inductive
current can produce large voltage spikes (the reason for flyback diodes and
snubbers).

\begin{exambox}
Remember the natural state: capacitor voltage and inductor current. That one
sentence tells you which quantity is continuous and which differential equation
to write.
\end{exambox}

\section{Worked Example}
A \(\SI{10}{\micro\henry}\) inductor is in series with a \(\SI{50}{\ohm}\)
resistor. The time constant and pole are
\[
\tau=\frac{L}{R}=\frac{10\times10^{-6}}{50}=\SI{2.0e-7}{\second},
\qquad
s_{\text{pole}}=-\frac{R}{L}=-\SI{5e6}{\per\second},
\]
so the corner frequency is \(f_c=1/(2\pi\tau)\approx\SI{796}{\kilo\hertz}\). At
\SI{14}{\mega\hertz}---well above the corner---the reactance
\(X_L=2\pi(14\times10^6)(10\times10^{-6})\approx\SI{880}{\ohm}\) dominates the
\SI{50}{\ohm} resistor, so the inductor's impedance sets the behavior, exactly as
the high-pass Bode plot predicts a decade past the corner.
